\subsection{初始研究目标}\label{route:problem}

最初的目标是改进一般 $3\times3$ 矩阵乘法的精确双线性复杂度：构造至多需要 22 个乘积的算法，或证明已知 23 个乘积的上界最优。目标针对双线性模型中的精确恒等式。在该模型中，每个计入次数的乘积，都是一个输入的线性形式与另一个输入的线性形式相乘。

Strassen 的七乘积 $2\times2$ 算法展示了双线性重组的力量\cite{strassen1969gaussian}。Laderman 给出了 $3\times3$ 的 23 乘积构造\cite{laderman1976noncommutative}，Bläser 发展了小尺寸情形的下界\cite{blaser2003}。基于 SAT 的算法发现与 AlphaTensor 表明，机器搜索能够发现精确乘法方案\cite{heule2019waysa,alphatensor2022}。新候选必须满足全部系数恒等式；下界则必须覆盖所有允许的分解。

两条大方向并行发展：一条探索有限域限制、组合支撑与证书；另一条探索特征零形变、秩一张成表示与对称性。它们不是固定方法之间的竞赛，而是可以相互提供例子、约束与交叉检查。

图\ref{fig:research-route}概括了这些联系。Qiushi Engine 在两个方向中发展了例子、表述和检验；它们的数学结果决定了哪些问题仍值得推进。

\subsection{张量约定与验证对照}\label{route:controls}

Qiushi Engine 首先建立统一的数学约定：按行优先将矩阵展平，从乘法规则生成张量，并对照全部 729 个张量元素检查已知方案。转置、循环置换、商映射和基变换都必须保持同一约定。

这不只是实现细节。表示线性泛函的矩阵与表示向量的矩阵，按互为对偶的方式变换；商类一般也不等于某个不变坐标补空间中的向量。混淆任意一对对象，都会改变编码中的数学问题，即使数组维度看起来仍然合理。

已知精确分解作为正对照；修改过的候选与故意不一致的实例检验拒绝机制。这些对照可以迅速证伪一个编码，但仅通过对照并不能证明编码覆盖每个合法分解。因此，最终证明还需要显式给出从假设分解到整数约束、再从整数约束到已检查公式的映射。

\subsection{形变、配对缺陷与数值证据的边界}\label{route:deformation}

一种自然的构造路线，是从已知秩 23 方案附近出发。研究考察了切方向、二阶相容性，以及趋向秩亏配对的延拓，目的是识别一个可以删除、同时仍保持目标张量的冗余乘积。

数值实验首先揭示了紧性问题：残差不断减小，并不保证候选趋近一个有限系数的解。若干看似有前景的轨迹在降低残差的同时，因子幅度不断增大。大项相互抵消，而非收敛至有限、精确的秩 22 分解。正则化揭示的是残差与幅度之间的权衡，而非精确解。所观察到的权衡明确了这些轨迹的失败机制；精确解空间中的其他分支仍有待研究。

随后，研究修订了表示方式。配对缺陷被表述为关联条件；是否张满所需空间、如何处理例外层，均与数值条件性分开考察。秩一张成表述考察：九维切片空间是否包含在由指定数目的秩一矩阵生成的张成空间中。补全坐标图试图更直接地实现这一几何约束。

过程中纠正了若干过度推论。满秩配对处的切空间计算，并不自动等于秩亏簇上的切锥计算；参数计数不是存在性定理；局部数值失败，即使发生于许多已知方案附近，也不能控制不连通分支或例外分支。这条路线没有产出特征零的秩 22 构造，也没有给出整体的秩 23 最优性证明。

\subsection{商核心与二元域上的问题重述}\label{route:cores}

Wang 的认证下界框架提供了另一入口：研究从因子空间移除若干方向之后剩下什么\cite{wangautomated}。下界 20 的证书被视为固定数学输入，并在全程使用的张量约定下重建其商语义。

$E_{11}$ 核心提供了一个有用的双向研究问题。移除这一第一因子方向后，张量形状为 $8\times9\times9$。其秩 19 构造可以嵌入完整张量，再加上三个秩一乘积恢复被删去的切片。因此，一个完整核心见证将给出 $\FF$ 上完整张量的至多秩 22 构造。

反方向则更为微妙。排除某个核心族，并不自动排除所有完整的 20 项分解。下界论证必须说明归约到该族的依据、对称性规范化、投影后的重数，以及对全部剩余分支的覆盖。研究反复回到这些待证明条件，而没有将一个困难子问题等同于整个定理。

核心有 255 个非零第一因子方向，完整的占据约束体系涉及 417,199 个子空间。这些都是有限、精确的对象，但可能的支撑仍构成巨大的组合空间。将目标重述于 $\FF$ 上，使搜索全程都可进行精确检验；识别有用结构仍然是核心困难。

\subsection{占据约束与因子补全}\label{route:completion}

对于长度为 $n$ 的分解，一个经过认证的商下界 $L(W)$ 意味着
\[
 m(W)\le n-L(W).
\]
该不等式记录有多少第一因子可以占据一个子空间，并有意舍去第二、第三因子的信息。因此，即使一个支撑满足全部占据不等式，也仍可能无法实现为分解。

Qiushi Engine 考察了固定第一因子的 SAT 表述、由收缩得到的共享因子条件，以及从更小商张量提升回来的问题。研究逐渐明确了三个层级：
\[
 \begin{gathered}
 \text{合法的投影支撑}\\
 \downarrow\\
 \text{到完整因子空间的相容提升}\\
 \downarrow\\
 \text{完整双线性分解}.
 \end{gathered}
\]
每个箭头都需要额外条件。较早层级的解，不是较晚层级的见证。

剩余容量搜索和根据违反情况学习的约束行，缩减了部分实例；但仅在选定前缀后才有效的约束，不能脱离其条件假设作全局复用。商空间提升比特的表述，也要求约束行导出器与分支求解器使用同一坐标映射。一次映射不一致曾产生表面上的排除结论，在重建语义后该结论被撤回。

重复投影方向作为独立分支处理，互异方向分支则仍然困难。最终的完整张量证明不需要确定整个 $E_{11}$ 核心的精确秩。

\subsection{受对称性限制的构造及其适用范围}\label{route:symmetry}

特征零方向探索了张量的内、外对称性、循环分解、转置不变族和混合对称类型。外幂与对称分量、剩余三次式、交换子或 Pfaffian 检验，以及秩一张成条件，提供了比无约束 Brent 方程更小的代数描述。

这条路线产生了局部障碍与精确检验工具。另一条同样重要的认识是：针对选定标架、种子轨道或对称性骨架得到的结果，不能提升为无约束的秩下界。有限域归约必须保留其底域范围；标架变换可能保持一个外幂分量，却改变下界检验实际使用的剩余对称分量。

若干表面上的捷径没有经受住检查。一项直和秩推论范围过大；某些检验在错误坐标中使用了对偶作用；对转置族提出的一个限制，与经过正确变换的已知构造比较后被否定。修正后的计算区分了完整张量的不变量与某个选定表示的性质。

这些探索为对称性构造提供了局部代数描述和检验方法。二元域证明最终沿另一条路线完成：使用适用于每个假设的 20 项分解的约束，而非选定对称族的约束。

\subsection{松弛与信息损失}\label{route:relaxation}

困难的互异支撑核心分支，促使研究考察编码理论、线性、成对相关和半定松弛。轨道平均的点对变量与矩矩阵压缩了大型带标签系统。精确的平方不等式有时可以将数值提示转化为可复用、可检查的条件。

这些实验澄清了证据的层级。分数占据向量不是整数支撑；整数点对计数型未必来自带标签点集；一个数值上半正定的矩阵本身不能提供所需因子分解。反之，病态的锥优化求解也不能证明精确不可行。

发现缺失的三角约束、被强迫的核、归一化因子或不同的约束行族后，某些中间结论必须缩小范围。能够分离一个已保存分数点的切平面，并不排除全部可行支撑。持续形成的认识是：每种压缩都必须说明保留了什么信息，还有什么需要重建。

这些松弛没有给出对整个核心的排除。因此，研究保留其中有用的约束，但没有将不断扩大松弛视为唯一的下一步。

\subsection{完整张量的秩约束}\label{route:full}

分裂展平将完整乘法张量重组为 $27\times27$ 置换矩阵 $P$。当另两个因子非零时，简单张量 $A_t\otimes B_t\otimes C_t$ 变成秩为 $\rk(A_t)$ 的矩阵 $M_t$。因此
\[
 27=\rk(P)\le\sum_{t=1}^{n}\rk(A_t).
\]
在 $n=20$ 时，相对二十个秩一第一因子的超额总和必须至少为七，故至少四个第一因子的矩阵秩至少为二。

这一界提供了仅靠第一因子占据约束未能鲜明表达的信息，也提示了一个更精简的问题：哪些低维商下界的加强，能够充分约束高秩因子对，从而确定其几何结构？

最终证明需要八个二维商轨道，其中每个平面的三个非零矩阵均至少秩二。将它们的商下界从 18 提高到 19，意味着在一个 20 项分解中，两个高秩第一因子不能张成这样的平面。因此，它们的和必须秩一。

\subsection{仿射几何与被强迫的秩型}\label{route:geometry}

依据经典矩阵几何\cite{hua1945geometries,wan1996geometry}，两两之差秩一的矩阵族包含于共同的仿射行陪集或列陪集。识别出这一结构后，成对的秩条件转化为一个可以显式分析占据情况的小型仿射构型。

在相关的规范化陪集中，八个第一行标签构成 $\FF^3$。相应矩阵有四个秩二、四个秩三。经过认证的仿射平面容量限制，使任意四点仿射平面上至多能选择三个点。$\FF^3$ 的每个五点子集都包含一个这样的平面，所以高秩因子至多有四个。

展平又要求至少四个。在四个高秩因子中，四个秩三点本身会形成被禁止的仿射平面。为了提供所需秩超额，唯一剩余可能是一个秩二、三个秩三。完整秩型为
\[
 (n_1,n_2,n_3)=(16,1,3),\qquad
 n_1+2n_2+3n_3=27.
\]
这一等式成为组合几何与代数刚性之间的关键连接。

\subsection{饱和与跨因子刚性}\label{route:saturation}

对可逆矩阵 $P=\sum_tM_t$，等式 $\sum_t\rk(M_t)=\rk(P)$ 使列空间构成直和。将向量按该直和分解，得到
\[
 M_tP^{-1}M_s=\delta_{ts}M_t.
\]
对于乘法张量，直接作指标收缩得到
\[
 M(A,B,C)P^{-1}M(A',B',C')
 =M\bigl(A(B'C^T)A',B,C'\bigr).
\]
当 $A_t$ 可逆时，对角恒等式强迫 $B_tC_t^T=A_t^{-1}$，所以 $B_t,C_t$ 均可逆。若两个第一因子 $A_t,A_s$ 都可逆，非对角恒等式要求 $A_tB_sC_t^TA_s=0$，而可逆矩阵的乘积不可能为零。

因此，一个饱和分解至多允许一个可逆第一因子；仿射秩型却要求三个。这正是证明下界 21 的矛盾。

Qiushi Engine 改变了最后的证明任务：一个可能很大的末端枚举变成了精确的跨因子恒等式。一个不必要的额外商轨道也从前提中删除。最终论证保留八个关键有限加强结果，并从依赖链中移除了更大的末端搜索。

\subsection{有限前提的认证}\label{route:certificates}

每个关键商空间有 127 个非零方向和 29,210 个非零真子空间。目标长度为 18 时，占据变量是非负整数，且总和被约束为 18；容量来自已验证的上游商下界表。

研究发展了两种编码。有限布尔副本将每个整数重数表示为若干二元变量之和；单调一元计数表示重数是否达到各个可能层级，后者直接呈现计数语义。当商允许重数二时，两种编码都必须允许它；将重数替换为互异支撑的比特变量，会改变数学问题。

此前一版基数编码意外复用了辅助变量，产生的是公式错误导致的矛盾，而不是对张量分解的排除。相应的表面下界结论已撤回。因此，原始 SAT 验证必须覆盖两部分：
\[
 \text{数学编码正确}
 \qquad+\qquad
 \text{公式不可满足性的证明正确}.
\]
DRAT 确立第二部分；显式语义重建、正对照和独立编码器检验第一部分。

最终八个子句证明全部保留，其公式可以重新生成并逐字节比较。源证书可以重新验证，其轨道代表可以独立展开为包含全部 8,283,458 个子空间的表。结构性论证随后从这些有限前提导出矛盾。

配套 Lean 开发在实际商空间对象上通过整数分支证书证明有限前提，并将它们连接到结构性矛盾。它还检查轨道分类、占据系统、校准对照与布尔编码语义。这些形式化对应结果保留了计算研究中的数学认识，同时使完整定理成为经过内核检查的依赖链（第\ref{sec:lean}节）。

\subsection{纠错及其方法论影响}\label{route:corrections}

\begin{longtable}{@{}L{3.3cm}L{5.1cm}L{5.8cm}@{}}
\toprule
问题&表面推论为何不成立&保留下来的科学认识\\
\midrule
\endfirsthead
\toprule 问题&表面推论为何不成立&保留下来的科学认识\\\midrule
\endhead
\bottomrule\endfoot
接近零的数值残差&系数可能在发散的同时相互抵消&要求有限系数的精确恒等式；区分已观察到的退化与不可能性\\
辅助变量冲突&布尔公式因编码原因自相矛盾&撤回结论；验证变量分配与语义正向映射\\
查找表解码&按错误的比特约定读取了压缩子空间键&重建规范行空间，并比较完整展开表\\
商空间坐标&投影与提升使用不同的坐标映射&明确保留商映射、补空间及诱导作用\\
局部对称性排除&种子、轨道或标架未覆盖全部分解&陈述精确假设，保留未覆盖分支\\
松弛的可行性&分数矩或平均矩未必实现为离散支撑&区分连续、整数、带标签及因子完整的对象\\
乘积指标&初等公式与实际置换收缩不一致&独立构造置换矩阵，检查初等乘积\\
不必要的有限前提&一个额外因子对轨道已经具有秩一差&从最短证明中移除，而非无必要地证明更强命题\\
\end{longtable}

纠错的价值体现在后续研究方法的改变上。基数错误促成了显式赋值映射与独立计数编码；坐标不一致促成了从张量及其映射出发的重建；发现某个有限前提不必要，促成了更短的论证。每次保留下来的知识，都比“认真检查”更具体：它明确了后续计算必须保持的数学关系。

\subsection{结果、可复用方法与开放问题}\label{route:outcome}

已确立的结果是
\[
 R_{\FF}(\Tmmm)\ge21.
\]
最终证明结合经过认证的商下界、仿射秩一几何、秩可加的分裂展平，以及张量特有的乘积恒等式。它还说明，在某个因子位置上零超额的秩 22 分解，至多有一个该位置的因子可逆；若该位置的全部 22 个因子均非零，剩余秩型为 $(17,5,0)$ 与 $(18,3,1)$。允许零因子时，仍保留归约到零超额 21 项分解的替代，见第\ref{sec:outlook}节。

这一下界及上述结构性推论均已有完整的 Lean 证明，其假设见 \qpath{formalization/COVERAGE.md}。主定理不再保留任何有限下界假设。

二元域上的精确秩仍位于 $[21,23]$。构造秩 22 算法、控制正超额分解，以及将机制扩展到其他域，是自然的后续问题。商张量和对称性材料在完整下界证明之外，也保留了通向这些问题的路线。

Meta-Trace 记录了第二项产出：研究方法如何随理解的积累而发展。Qiushi Engine 构造表示，用反例与计算识别其边界，并将有用区别带入后续推理。支撑与补全之间的区别，转化为如何恢复因子耦合的问题；被强迫的取等关系给出了答案。

将这一发展过程与数学材料共同记录，不仅使已接受的证明能够复验。读者还可以检查实验背后的假设，理解结果为何改变路线，改造中间方法，或在明确既往障碍的基础上重新研究开放问题。证明确立结果，Meta-Trace 则保留可供后续研究继续发展的推理历史。
